\section{Проведённые эксперименты и исследование модели}
\subsection{Цель и схема исследования}
Исследование выполнено на текущем генераторе \code{calendar-v7} 6 ноября 2026 года. Цель --- сравнить актуализацию календаря, отказы и завершённые мероприятия при трёх способах обработки конфликтов и двух интенсивностях поступления заявок. Исходные данные получены запуском существующего консольного средства проекта, без изменения алгоритма ради результатов.
\begin{lstlisting}
build/reorganized/modeling_experiments \
  --output documentation/experiments/results-v7 \
  --repetitions 20
\end{lstlisting}

План включает шесть сочетаний: \code{Reject}, \code{Suggest} и \code{AutoReschedule} при интенсивностях 6 и 12 случайных заявок на фирму за рабочий день. На сочетание выполнено 20 опытов с seed 42--61, всего 120 запусков. Для пакетного \code{Suggest} задан явный оператор: первый предложенный вариант принимается после каждого шага. В обычном интерфейсе такой выбор делает пользователь.

\begin{longtable}{L{0.47\linewidth}L{0.45\linewidth}}
\toprule Параметр & Значение исследования \\\midrule\endhead
Начало и период & 12.10.2026, 14 дней \\
Отделы и штат & 6 отделов, по 8 сотрудников, всего 48 \\
Помещения & 6 комнат отделов и конференц-зал \\
Рабочее время и шаг & Пн--Пт, 09:00--18:00; 30 минут \\
Исходный календарь & 60 планёрок и 2 совещания руководителей \\
Периодические семинары & Включены; период 7 дней, заблаговременность 2 дня, 14:00, 90 минут \\
Случайные заявки и изменения & 6/12 заявок; 1 попытка изменения за рабочий день \\
Участники случайного плана & 2--8; без повторов \\
Дальний срок планирования & 1--5 рабочих дней; отдельная ветвь текущего дня \\
Веса важности & 0.5; 0.35; 0.15 \\
Напоминания & За 1/2/4 календарных дня, повторы каждые 24 часа \\
Поиск переносов & 5 рабочих дней, максимум 3 варианта \\
\bottomrule
\end{longtable}

\subsection{Статистические факторы и показатели}
За рабочую минуту число новых заявок распределено по Пуассону с параметром $\lambda/M$, где $M=540$ --- число минут рабочего дня, $\lambda$ --- заданная дневная интенсивность. Независимо формируются попытки изменения расписания. Тип, важность и допустимые участники выбираются по дискретным и равномерным распределениям; начало выбирается среди допустимых минут рабочего дня.

Основные показатели: созданные и завершённые мероприятия, конфликтные заявки, окончательно отказанные попытки, имитируемые доставки и доля занятости помещений. Будущие планы, выходящие за горизонт, могут быть созданы, но ещё не завершены к его концу. Поэтому число созданных мероприятий не равно числу завершённых.

Для $n=20$ наблюдений показателя $x$ вычисляются
\[
 \bar x=\frac1n\sum_{i=1}^{n}x_i,\qquad
 s_x=\sqrt{\frac{\sum_{i=1}^{n}(x_i-\bar x)^2}{n-1}}.
\]
Доля занятости --- сумма рабочих секунд фактически начавшихся событий, обрезанных прошедшим временем, делённая на прошедший рабочий фонд всех помещений. Будущие бронирования не включаются.

\subsection{Фактические результаты}
В таблице приведены средние значения и выборочные стандартные отклонения в формате $\bar x\pm s_x$. Обозначения: R --- отклонение, S --- предложения с выбором первого варианта после шага, A --- немедленный автоматический перенос; $\lambda$ --- дневная интенсивность. Подробные наблюдения доступны в JSON и CSV набора \file{results-v7}.
\begin{table}[H]\centering\small
\setlength{\tabcolsep}{3pt}
\begin{tabular}{crrrr}
\toprule
Режим & $\lambda$ & Завершено & Отказано & Занятость, \% \\
\midrule
R & 6 & $93.05\pm3.71$ & $20.00\pm4.12$ & $9.46\pm0.52$ \\
S & 6 & $107.60\pm5.78$ & $0.05\pm0.22$ & $11.32\pm0.82$ \\
A & 6 & $107.55\pm5.77$ & $0.05\pm0.22$ & $11.31\pm0.82$ \\
R & 12 & $119.05\pm5.03$ & $50.80\pm8.84$ & $12.80\pm0.91$ \\
S & 12 & $153.45\pm10.76$ & $0.25\pm0.64$ & $17.42\pm1.81$ \\
A & 12 & $153.25\pm10.60$ & $0.20\pm0.52$ & $17.40\pm1.77$ \\
\bottomrule
\end{tabular}
\caption{Результаты 120 прогонов calendar-v7: по 20 наблюдений на сочетание.}
\end{table}

\paragraph{Интенсивность 6.} При отклонении конфликтов завершено в среднем 93.05 мероприятий; при автоматическом переносе --- 107.55. Разница составляет 14.50 мероприятия. Среднее число отказанных попыток соответственно 20.00 и 0.05; имитируемых доставок --- 1229.70 и 1364.05.

\paragraph{Интенсивность 12.} При отклонении конфликтов завершено в среднем 119.05 мероприятий; при автоматическом переносе --- 153.25. Разница составляет 34.20 мероприятия. Среднее число отказанных попыток соответственно 50.80 и 0.20; имитируемых доставок --- 1414.85 и 1755.25.

\paragraph{Наблюдение по стратегии предложений.} Результаты Suggest с выбором первого варианта близки к немедленному автоматическому переносу. Это согласуется с заданным программным правилом решения, но не доказывает равенство стратегий при другом поведении оператора. Большее число конфликтных заявок при переносах не означает ухудшение проверки: более заполненный календарь создаёт дополнительные пересечения, которые могут завершаться успешным переносом.


\subsection{Ограничения интерпретации}
\begin{itemize}
\item Веса и длительности демонстрационные, эмпирическая калибровка по данным реальной фирмы не проводилась.
\item При выборе Suggest возникает новая проверяемая попытка. Поэтому суммарное число заявок между стратегиями не является одним и тем же фиксированным внешним потоком.
\item Одинаковый seed не обеспечивает одинаковую последовательность последующих событий между стратегиями: результаты операций меняют состояние и ветвление генерации.
\item Эксперимент оценивает модель с автоматическим оператором Suggest, а не поведение реального пользователя.
\item Значимость различий, причинный эффект и универсальное превосходство стратегии не устанавливались. Выводы относятся к указанным параметрам и выполненной серии.
\end{itemize}

Набор v4 в \file{documentation/experiments/results} является историческим. Его наблюдения не использованы в таблицах текущего отчёта. Скрипт \file{make_results.py} перепроверяет балансы, средние и стандартные отклонения v7 и фиксирует источник и контрольную сумму в \file{data_manifest.json}.
